\chapter{Sous-variétés de \texorpdfstring{$\R^n$}{Rn}}
\label{ch:b3:submanifolds}

Sphères, tores, groupes de rotations~: les habitats naturels de
la géométrie et de la mécanique ne sont pas des espaces
vectoriels mais des \emph{ensembles courbes qui paraissent
plats de près}. Ce chapitre donne un sens précis à cette phrase
--- les sous-variétés de $\R^n$ --- et le calcul pour travailler
dessus. Le fondement est le théorème d'inversion locale, prouvé
ici par le point fixe de Banach~; tout le reste est un
changement de coordonnées~: les quatre descriptions équivalentes
d'une sous-variété (redressement local, ensembles de niveau,
graphes, paramétrisations), les \omterm{def:b3:submanifolds:tangent}{espaces tangents}, et
l'optimisation sous contrainte par les \omterm{thm:b3:submanifolds:lagrange}{multiplicateurs de
Lagrange} --- qui, en démonstration d'adieu, reprouve le théorème
spectral pour les matrices symétriques en trois lignes de
géométrie. Le problème de week-end construit le groupe de
rotations $SO(3)$ et son double revêtement quaternionique~:
l'algèbre ($Q_8$ du~\cref{ch:b3:groups} devenu adulte)
rencontrant la géométrie.

\section{Le théorème d'inversion locale}

\begin{theorem}[Théorème d'inversion locale]
\label{thm:b3:submanifolds:ift}
Soit $U \subseteq \R^n$ \omterm{def:b3:topology:topology}{ouvert}, $f \colon U \to \R^n$ de
classe $\mathcal C^1$, et $a \in U$ avec $Df(a)$ inversible.
Alors il existe des \omterm{def:b3:topology:topology}{ouverts} $V \ni a$, $W \ni f(a)$ tels que
$f \colon V \to W$ est une bijection d'inverse $\mathcal C^1$,
et\index{théorème d'inversion locale}
\[
D(f^{-1})(y) = \bigl(Df(f^{-1}(y))\bigr)^{-1}
\qquad (y \in W).
\]
Si $f$ est $\mathcal C^k$, il en va de même de $f^{-1}$.
\end{theorem}

\begin{proof}
Normalisons~: en remplaçant $f$ par $x \mapsto Df(a)^{-1}\bigl(f(a +
x) - f(a)\bigr)$, on peut supposer $a = 0$, $f(0) = 0$, $Df(0) =
I$ (l'énoncé général suit en composant avec les bijections
affines). Écrivons $f(x) = x + g(x)$~: $Dg(0) = 0$, et par
continuité de $Dg$ choisissons $r > 0$ avec
$\vertiii{Dg(x)} \leq \frac12$ sur $\bar B(0, r)$~; l'inégalité
des accroissements finis donne $\norm{g(x) - g(x')} \leq \frac12
\norm{x - x'}$ là.

\emph{Bijectivité sur un \omterm{def:b3:topology:topology}{voisinage}.} Pour $y \in B(0,
\frac r2)$, résoudre $f(x) = y$ revient à trouver un point fixe
de $\Phi_y(x) = y - g(x)$~; $\Phi_y$ envoie $\bar B(0, r)$ dans
elle-même ($\norm{\Phi_y(x)} \leq \norm y + \frac12\norm x \leq
r$) et est $\frac12$-lipschitzienne~: Banach
(\cref{thm:b3:complete:banach}) fournit une unique solution
$x = \varphi(y) \in \bar B(0,r)$. De plus $f$ est injective
sur $\bar B(0,r)$~:
\[
\norm{f(x) - f(x')} \geq \norm{x - x'} - \norm{g(x) - g(x')}
\geq \tfrac12\norm{x - x'} .
\tag{$*$}
\]
Posons $W = B(0, \frac r2)$ et $V = f^{-1}(W)\cap B(0, r)$~:
\omterm{def:b3:topology:topology}{ouvert} (continuité), avec $f \colon V \to W$ bijective.

\emph{Continuité et différentiabilité de l'inverse.} $(*)$
dit que $\varphi = f^{-1}$ est $2$-lipschitzienne. Fixons $y_0 = f(x_0)
\in W$~; l'inversibilité de $A = Df(x_0)$ (sa distance à $I$
est $\leq \frac12$~: Neumann,
\cref{prop:b3:banach:neumann}) et la différentiabilité de $f$
donnent, pour $y = f(x)$ près de $y_0$~:
\[
\varphi(y) - \varphi(y_0) - A^{-1}(y - y_0)
= -A^{-1}\bigl(f(x) - f(x_0) - A(x - x_0)\bigr)
= -A^{-1}\,o(\norm{x - x_0}) = o(\norm{y - y_0}),
\]
en utilisant $(*)$ pour convertir $\norm{x - x_0} \leq 2\norm{y - y_0}$~:
$\varphi$ est différentiable en $y_0$ avec la différentielle
inverse. Continuité de $y \mapsto D\varphi(y) =
Df(\varphi(y))^{-1}$~: composition d'\omterm{def:b3:topology:continuity}{applications continues}
(l'inversion est \omterm{def:b3:topology:continuity}{continue},
\cref{prop:b3:banach:neumann})~: $\varphi \in \mathcal C^1$~;
le bootstrap de la même formule donne $\mathcal C^k$.
\end{proof}

\begin{theorem}[Théorème des fonctions implicites]
\label{thm:b3:submanifolds:implicit}
Soit $F \colon U \subseteq \R^p\times\R^q \to \R^q$ de classe
$\mathcal C^1$ près de $(a, b)$, $F(a,b) = 0$, et supposons la
différentielle partielle $D_yF(a,b) \in \mathcal L(\R^q)$
inversible. Alors il existe des \omterm{def:b3:topology:topology}{voisinages} $A \ni a$, $B \ni b$
et une application $\mathcal C^1$ $\psi \colon A \to B$ avec
\[
\bigl\{(x, y)\in A\times B : F(x,y) = 0\bigr\}
= \{(x, \psi(x)) : x \in A\},
\]
et $D\psi(x) = -D_yF(x, \psi(x))^{-1}\,D_xF(x,
\psi(x))$.\index{théorème des fonctions implicites}
\end{theorem}

\begin{proof}
Appliquer le~\cref{thm:b3:submanifolds:ift} à $\Theta(x, y) = (x,
F(x,y))$~: sa différentielle en $(a,b)$, triangulaire par blocs
avec blocs diagonaux inversibles $I$ et $D_yF$, est inversible.
L'inverse local est de la forme $\Theta^{-1}(x, z) = (x, h(x,
z))$~; posons $\psi(x) = h(x, 0)$~: alors $F(x, y) = 0$ ssi
$\Theta(x,y) = (x, 0)$ ssi $y = \psi(x)$, localement. La
formule~: dériver $F(x, \psi(x)) = 0$ par la règle de la
chaîne.
\end{proof}

\section{Sous-variétés~: quatre définitions}

\begin{theorem}[Caractérisations équivalentes]
\label{thm:b3:submanifolds:characterizations}
Soit $M \subseteq \R^n$, $d \in \{0, \dots, n\}$, et $k \geq
1$. Les assertions suivantes sont équivalentes, pour chaque
point $a \in M$ (et $M$ est une \emph{sous-variété de dimension
$d$ de classe $\mathcal C^k$}\index{sous-variété} si elles
valent en tout $a \in M$)~:
\begin{enumerate}
  \item (\emph{Redressement}) Il existe un difféomorphisme
        $\mathcal C^k$ $\Phi$ d'un \omterm{def:b3:topology:topology}{ouvert} $\Omega \ni a$ sur
        un \omterm{def:b3:topology:topology}{ouvert} $\Omega' \subseteq \R^n$ avec
        \[
        \Phi(M\cap\Omega) = \Omega' \cap
        \bigl(\R^d\times\{0\}\bigr).
        \]
  \item (\emph{Ensemble de niveau}) Il existe une
        \emph{submersion} $\mathcal C^k$ $F \colon \Omega \to
        \R^{n-d}$ (c'est-à-dire $DF(x)$ surjective) sur un
        \omterm{def:b3:topology:topology}{ouvert} $\Omega \ni a$ avec $M\cap\Omega = F^{-1}(0)$.
  \item (\emph{Graphe}) À permutation des coordonnées près,
        $M$ est localement le graphe d'une application
        $\mathcal C^k$ $\psi \colon A \subseteq \R^d \to
        \R^{n-d}$.
  \item (\emph{Paramétrisation}) Il existe une
        \emph{immersion} $\mathcal C^k$ $\varphi \colon A
        \subseteq \R^d \to \R^n$ ($D\varphi(u)$ injective)
        avec $A$ \omterm{def:b3:topology:topology}{ouvert}, $\varphi$ un \omterm{def:b3:topology:continuity}{homéomorphisme} de $A$
        sur $M \cap \Omega$ pour un \omterm{def:b3:topology:topology}{ouvert} $\Omega \ni a$.
\end{enumerate}
\end{theorem}

\begin{proof}
(1)$\Rightarrow$(2)~: $F = (\Phi_{d+1}, \dots, \Phi_n)$
(dernières coordonnées de $\Phi$)~: une submersion ($D\Phi$
inversible).
(2)$\Rightarrow$(3)~: $DF(a)$ surjective~: un mineur $q
\times q$ de la jacobienne est inversible ($q = n - d$)~;
après permutation des coordonnées, $D_yF(a)$ est inversible, et
le théorème des fonctions implicites
(\cref{thm:b3:submanifolds:implicit}) exprime $M$ localement
comme un graphe $y = \psi(x)$.
(3)$\Rightarrow$(4)~: $\varphi(x) = (x, \psi(x))$~: une
immersion (différentielle $\bigl(\begin{smallmatrix}I\\
D\psi\end{smallmatrix}\bigr)$ injective), un \omterm{def:b3:topology:continuity}{homéomorphisme} sur
le graphe (inverse~: la projection, \omterm{def:b3:topology:continuity}{continue}).
(4)$\Rightarrow$(1)~: soit $\varphi(u_0) = a$~; compléter
$\operatorname{im}D\varphi(u_0)$ par un supplémentaire $E$
($\dim E = n - d$) et définir $\Theta(u, v) = \varphi(u) + v$
sur $A\times E$~: $D\Theta(u_0, 0)$ est bijective (l'image
contient $\operatorname{im}D\varphi(u_0)$ et $E$), donc
$\Theta$ est un difféomorphisme local
(\cref{thm:b3:submanifolds:ift})~; son inverse $\Phi$
redresse~: près de $a$, les points de $M$ sont exactement les
$\varphi(u) = \Theta(u, 0)$ --- pour cela, l'hypothèse
d'\omterm{def:b3:topology:continuity}{homéomorphisme} dans (4) garantit que $M\cap\Omega$, pour
$\Omega$ petit, ne contient pas d'autres feuilles
($\varphi(u') + v = m \in M$ proche de $a$ avec $v \neq 0$
petit doit être exclu~: $m = \varphi(u'')$ pour un $u''$ près
de $u_0$ par la propriété d'\omterm{def:b3:topology:continuity}{homéomorphisme}, et l'injectivité
locale de $\Theta$ force $v = 0$). Alors $\Phi(M\cap\Omega) =
(A\times\{0\})\cap\Phi(\Omega)$ quitte à rétrécir.
\end{proof}

\begin{example}\label{ex:b3:submanifolds:examples}
La sphère $S^{n-1} = \{\norm x^2 = 1\}$~: ensemble de niveau
de la submersion $F(x) = \norm x_2^2 - 1$ sur
$\R^n\setminus\{0\}$ ($DF(x) = 2x^{\mathsf T} \ne 0$)~: une
sous-variété $\mathcal C^\infty$ de dimension $n - 1$. Le tore
dans $\R^3$~: ensemble de niveau de $\bigl(\sqrt{x^2 + y^2} -
R\bigr)^2 + z^2 - r^2$ ($0 < r < R$). Le cône $\{x^2 + y^2 =
z^2\}$ n'est \emph{pas} une sous-variété en $0$
(\cref{exo:b3:submanifolds:1}). Groupes de matrices~: $SL_n$ et
$O_n$ sont des sous-variétés de $M_n(\R)$
(\cref{exo:b3:submanifolds:5,exo:b3:submanifolds:6}) --- le
point de départ de la théorie de Lie.
\end{example}

\section{Espaces tangents}

\begin{definition}\label{def:b3:submanifolds:tangent}
Soit $M$ une $d$-sous-variété et $a \in M$. L'\emph{espace
tangent}\index{espace tangent} $T_aM$ est l'ensemble des
vecteurs vitesse $\gamma'(0)$ des courbes $\mathcal C^1$
$\gamma \colon \intoo{-\varepsilon}\varepsilon \to M$ avec
$\gamma(0) = a$.
\end{definition}

\begin{proposition}\label{prop:b3:submanifolds:tangent}
$T_aM$ est un sous-espace vectoriel de dimension $d$ de
$\R^n$, et~:
\begin{enumerate}
  \item si $M = F^{-1}(0)$ localement avec $F$ une submersion~:
        $T_aM = \ker DF(a)$~;
  \item si $M$ est paramétrée par l'immersion $\varphi$
        ($\varphi(u_0) = a$)~: $T_aM =
        \operatorname{im}D\varphi(u_0)$.
\end{enumerate}
\end{proposition}

\begin{proof}
Les courbes dans $M$ vérifient $F(\gamma(t)) = 0$~; la règle
de la chaîne en $0$ donne $DF(a)\gamma'(0) = 0$~: $T_aM
\subseteq \ker DF(a)$. Réciproquement, le redressement
(\cref{thm:b3:submanifolds:characterizations}(1)) transporte
les droites de $\R^d\times\{0\}$ en courbes de $M$~: tout
vecteur d'un sous-espace de dimension $d$ est réalisé~; en
comparant les dimensions ($\dim\ker DF(a) = n - (n - d) =
d$) on force l'égalité dans (1), et le même argument de
transport donne (2) ($D\varphi(u_0)$ appliqué aux droites de
$A$~; dimensions encore).
\end{proof}

\begin{theorem}[Multiplicateurs de Lagrange]
\label{thm:b3:submanifolds:lagrange}
Soit $M = F^{-1}(0)$ avec $F = (F_1, \dots, F_q) \colon
\Omega \to \R^q$ une submersion $\mathcal C^1$, et $f \colon
\Omega \to \R$ de classe $\mathcal C^1$. Si la restriction
$f\restriction_M$ a un extremum local en $a \in M$, alors il
existe d'uniques réels $\lambda_1, \dots, \lambda_q$
(\emph{\omterm{thm:b3:submanifolds:lagrange}{multiplicateurs de Lagrange}}) avec\index{multiplicateurs de Lagrange}
\[
\nabla f(a) = \lambda_1\nabla F_1(a) + \dots +
\lambda_q\nabla F_q(a) .
\]
\end{theorem}

\begin{proof}
Pour toute courbe $\gamma$ dans $M$ passant par $a$~: $t
\mapsto f(\gamma(t))$ a un extremum local en $0$, donc $0 =
\frac{\dd}{\dd t}f(\gamma(t))\big|_0 = \langle\nabla f(a),
\gamma'(0)\rangle$~: $\nabla f(a) \perp T_aM = \ker DF(a)$
(\cref{prop:b3:submanifolds:tangent}). Or $\ker DF(a)^\perp
= \operatorname{im}DF(a)^{\mathsf T}$ (l'identité
rang/orthogonalité de l'année~2, ou \cref{exo:b3:hilbert:8}
en dimension finie~: $(\ker T)^\perp = \operatorname{im}T^*$),
qui est engendré par les gradients $\nabla F_i(a)$ ---
indépendants, car $DF(a)$ est surjective~: les multiplicateurs
existent et sont uniques.
\end{proof}

\begin{example}[Le théorème spectral, géométriquement]
\label{ex:b3:submanifolds:rayleigh}
Soit $A$ une matrice réelle symétrique $n\times n$ et
maximisons $f(x) = \langle Ax, x\rangle$ sur la sphère
$S^{n-1}$ (compacte~: le maximum est atteint, en un certain
$v_1$). Lagrange avec $F(x) = \norm x^2 - 1$~: $\nabla f =
2Ax$ et $\nabla F = 2x$ donnent $Av_1 = \lambda_1v_1$ --- un
vecteur propre, avec $\lambda_1 = \max_{S^{n-1}}\langle Ax,
x\rangle$. Restreignons $A$ à $v_1^\perp$ (invariant~:
$\langle Av, v_1\rangle = \langle v, Av_1\rangle =
\lambda_1\langle v, v_1\rangle = 0$) et itérons~: une \omterm{def:b3:topology:basis}{base}
orthonormée de vecteurs propres. Le théorème spectral de
l'année~2, reprouvé par pure optimisation --- et l'ombre de
dimension infinie du même argument a prouvé
le~\cref{lem:b3:spectral:existence}.
\end{example}

\begin{method}\label{met:b3:submanifolds:method}
Pour prouver qu'un ensemble est une sous-variété~: l'exhiber
localement comme $F^{-1}(0)$ avec $DF$ surjective \emph{sur
l'ensemble} (la voie la plus courante), ou comme un graphe.
Pour calculer sa dimension et son \omterm{def:b3:submanifolds:tangent}{espace tangent}~: $d = n - q$
et $T_a = \ker DF(a)$. Pour optimiser dessus~: Lagrange ---
toujours vérifier d'abord la compacité (ou la coercivité),
pour qu'un extremum existe auquel le théorème puisse
s'appliquer, et se souvenir que l'équation des multiplicateurs
n'est que \emph{nécessaire}~: collecter tous les points
critiques, puis comparer les valeurs. Pour les groupes de
matrices, dériver les courbes en l'identité pour identifier
les \omterm{def:b3:submanifolds:tangent}{espaces tangents}.
\end{method}

\section{Exercices}

\begin{exercise}[$\star$]\label{exo:b3:submanifolds:1}
(a) Vérifier que les suivants sont des sous-variétés
$\mathcal C^\infty$ et donner leurs dimensions~: $S^{n-1}$~;
l'hyperboloïde $\{x^2 + y^2 - z^2 = 1\}$~; le tore de
l'\cref{ex:b3:submanifolds:examples}.
(b) Montrer que le cône $C = \{x^2 + y^2 = z^2\} \subseteq
\R^3$ n'est pas une $2$-sous-variété en $0$~: déterminer le
nombre de \omterm{def:b3:topology:components}{composantes connexes} de
$\bigl(C\setminus\{0\}\bigr)\cap B(0,\varepsilon)$, et
comparer avec le nombre qu'un redressement
(\cref{thm:b3:submanifolds:characterizations}(1)) forcerait
pour un plan privé d'un point.
\end{exercise}

\begin{exercise}[$\star$]\label{exo:b3:submanifolds:2}
Calculer les \omterm{def:b3:submanifolds:tangent}{espaces tangents}~:
(a) $T_aS^{n-1}$ pour tout $a$ (réponse~: $a^\perp$)~;
(b) le plan tangent au tore de
l'\cref{ex:b3:submanifolds:examples} en un point arbitraire
de l'équateur extérieur $\{z = 0,\ x^2 + y^2 = (R + r)^2\}$~;
(c) la droite tangente à l'hélice $\varphi(t) = (\cos t,
\sin t, t)$ en $\varphi(t_0)$, en vérifiant
le~\cref{prop:b3:submanifolds:tangent}(2).
\end{exercise}

\begin{exercise}[$\star\star$]\label{exo:b3:submanifolds:3}
Soit $f(x, y) = (x^2 - y^2,\ 2xy)$ (c'est-à-dire $z \mapsto
z^2$).
(a) En quels points le~\cref{thm:b3:submanifolds:ift}
s'applique-t-il~?
(b) Montrer que $f$ est localement mais non globalement
inversible sur $\R^2\setminus\{0\}$, et exhiber explicitement
les deux inverses locaux définis sur un \omterm{def:b3:topology:topology}{voisinage} de $(1, 0)$
(les deux branches de racine carrée).
(c) Même discussion pour l'application des \omterm{ex:b3:product:polar}{coordonnées
polaires} $(r, \theta) \mapsto (r\cos\theta, r\sin\theta)$.
\end{exercise}

\begin{exercise}[$\star\star$]\label{exo:b3:submanifolds:4}
(Folium) Soit $F(x,y) = x^3 + y^3 - 3xy$ et $\mathcal C =
F^{-1}(0)$.
(a) Montrer que près de tout point de $\mathcal C$ autre que
l'origine, $\mathcal C$ est une $1$-sous-variété, localement
un graphe en $x$ ou en $y$ (lequel, où~?).
(b) Calculer la droite tangente en $(\frac32, \frac32)$.
(c) Que se passe-t-il à l'origine~? (Deux branches se
croisent~: exhiber deux courbes $\mathcal C^1$ dans $\mathcal
C$ passant par $0$ à vitesses indépendantes, et conclure
qu'aucun redressement n'existe.)
\end{exercise}

\begin{exercise}[$\star\star$]\label{exo:b3:submanifolds:5}
Soit $F(M) = M^{\mathsf T}M$ de $M_n(\R)$ vers l'espace
$S_n$ des matrices symétriques.
(a) Montrer $DF(M)(H) = M^{\mathsf T}H + H^{\mathsf T}M$ et
que $DF(M)$ est surjective sur $S_n$ en tout $M \in O_n$
\emph{(étant donné $S \in S_n$, essayer $H = \frac12MS$)}.
(b) Conclure que $O_n = F^{-1}(I)$ est une sous-variété
$\mathcal C^\infty$ compacte de dimension
$\frac{n(n-1)}2$, avec $T_IO_n = \{H : H^{\mathsf T} = -H\}$,
les matrices antisymétriques.
(c) Montrer que $\eu^{tH} \in O_n$ pour tout $H$
antisymétrique~: les directions tangentes s'intègrent en
courbes dans le groupe.
\end{exercise}

\begin{exercise}[$\star\star$]\label{exo:b3:submanifolds:6}
(a) Montrer que $\det \colon M_n(\R) \to \R$ a pour
différentielle $D\det(M)(H) =
\operatorname{tr}\bigl(\operatorname{com}(M)
^{\mathsf T}H\bigr)$, non nulle en tout $M \in SL_n$.
(b) Conclure que $SL_n(\R)$ est une sous-variété de dimension
$n^2 - 1$ avec $T_ISL_n = \{H : \operatorname{tr}H = 0\}$.
(c) $GL_n(\R)$ est-il une sous-variété~? De quelle
dimension~?
\end{exercise}

\begin{exercise}[$\star\star$]\label{exo:b3:submanifolds:7}
Par les \omterm{thm:b3:submanifolds:lagrange}{multiplicateurs de Lagrange}~:
(a) trouver les extrema de $f(x,y) = xy$ sur le cercle $x^2 +
y^2 = 1$~;
(b) montrer que parmi tous les vecteurs de probabilité
$(p_1, \dots, p_n)$ (positifs, de somme $1$), l'entropie
$-\sum p_i\ln p_i$ est maximisée exactement en la
distribution uniforme~;
(c) trouver le point de l'ellipse $\{x^2/4 + y^2 = 1\}$ le
plus proche de $(1, 0)$, et vérifier l'équation des
multiplicateurs géométriquement (alignement des \omterm{def:b3:groups:normal}{normales}).
\end{exercise}

\begin{exercise}[$\star\star\star$]\label{exo:b3:submanifolds:8}
Rédiger pleinement l'\cref{ex:b3:submanifolds:rayleigh}~:
prouver par récurrence qu'une matrice réelle symétrique
admet une \omterm{def:b3:topology:basis}{base} orthonormée de vecteurs propres, avec
$\lambda_1 \geq \dots \geq \lambda_n$ les maxima contraints
successifs du quotient de Rayleigh. Puis déduire les formules
de Courant--Fischer de l'\cref{exo:b3:spectral:8} en
dimension finie directement de cette construction.
\end{exercise}

\begin{exercise}[$\star\star\star$]\label{exo:b3:submanifolds:9}
(Inégalité de Hadamard) Pour $M \in GL_n(\R)$ de colonnes
$c_1, \dots, c_n$~:
\[
\abs{\det M} \;\leq\; \prod_{i=1}^n\norm{c_i}_2 ,
\]
avec égalité ssi les colonnes sont orthogonales. \emph{(Se
ramener à des colonnes de norme $1$ par homothétie~;
maximiser $\det$ sur le produit compact de sphères
$(S^{n-1})^n$~; en un maximiseur, Lagrange dans chaque colonne
séparément donne $\nabla_{c_i}\det = \lambda_ic_i$, et
$\nabla_{c_i}\det$ est la $i$-ème colonne de
$\operatorname{com}(M)$~: en déduire que $M^{\mathsf T}M$
est diagonale, donc $= I$, donc $\det = \pm1$.)} Lecture
géométrique~: le volume d'un parallélépipède est au plus le
produit des longueurs de ses arêtes.
\end{exercise}

\begin{exercise}[$\star\star$]\label{exo:b3:submanifolds:10}
Près de lesquels de ses points le cercle $S^1$ est-il un
graphe $y = \psi(x)$~? Un graphe $x = \chi(y)$~? Vérifier la
caractérisation par graphe
(\cref{thm:b3:submanifolds:characterizations}(3))
explicitement en $(1, 0)$, et expliquer en une phrase
pourquoi \emph{une certaine} permutation de coordonnées
suffit toujours mais aucune n'est toujours adaptée.
\end{exercise}

\begin{exercise}[$\star\star$]\label{exo:b3:submanifolds:11}
(Le groupe orthogonal comme sous-variété, quantitativement)
(a) Montrer que $O_n = \{M : M^{\mathsf T}M = I\}$ est
compact~: borné (chaque colonne est un vecteur unitaire, donc
$\norm M \leq \sqrt n$ pour la norme euclidienne de
matrice) et fermé.
(b) Montrer que son \omterm{def:b3:submanifolds:tangent}{espace tangent} en $I$ est l'espace des
matrices antisymétriques, de dimension $\frac{n(n-1)}2$, et
en un $A \in O_n$ général~: $T_AO_n = \{AK : K^{\mathsf T} =
-K\}$.
(c) En déduire que l'application $t \mapsto A\exp(tK)$ est,
pour chaque $K$ antisymétrique, une courbe dans $O_n$ passant
par $A$ de vitesse $AK$ \emph{(vérifier $\exp(tK) \in O_n$
via $\exp(X)^{\mathsf T} = \exp(X^{\mathsf T})$ et
$\exp(-X)\exp(X) = I$)}~: tout vecteur tangent est réalisé
par une courbe explicite, sans théorème des fonctions
implicites.
\end{exercise}

\begin{exercise}[$\star\star$]\label{exo:b3:submanifolds:12}
(Points critiques de la distance) Soit $M \subseteq \R^n$ une
sous-variété et $p \notin M$. Montrer que si $x_0 \in M$
minimise la distance à $p$ (un tel point existe quand $M$ est
fermée et non vide --- pourquoi~?), alors
\[
p - x_0 \;\perp\; T_{x_0}M
\]
\emph{(dériver $t \mapsto \norm{\gamma(t) - p}^2$ le long des
courbes dans $M$)}. En déduire~: le point le plus proche sur
une sphère se trouve sur le rayon passant par le \omterm{ex:b3:groups:actions}{centre}~; et
utiliser la condition pour calculer la distance de $p =
(2, 0)$ à la parabole $y = x^2$ (se ramener à une cubique et
la résoudre numériquement à trois chiffres).
\end{exercise}

\section{Problème~: \texorpdfstring{$SO(3)$}{SO(3)} et les
quaternions}

\begin{problem}[{Problème de week-end --- rotations, le groupe
$S^3$, et le double revêtement}]\label{pb:b3:submanifolds:1}
Les quaternions $\mathbb H = \{t + x\mathrm i + y\mathrm j +
z\mathrm k\}$ --- l'algèbre dont le groupe des unités
contient le $Q_8$ du~\cref{pb:b3:groups:1} --- paramétrisent
les rotations tridimensionnelles en double~: l'application
«~conjuguer par un quaternion unitaire~» est un morphisme
surjectif $S^3 \to SO(3)$ de noyau $\{\pm1\}$. Nous
construisons tout. Rappel/définition~: la multiplication est
$\R$-bilinéaire avec $\mathrm i^2 = \mathrm j^2 = \mathrm
k^2 = \mathrm{ijk} = -1$~; le conjugué de $q = t +
x\mathrm i + y\mathrm j + z\mathrm k$ est $\bar q = t -
x\mathrm i - y\mathrm j - z\mathrm k$~; $N(q) = q\bar q =
t^2 + x^2 + y^2 + z^2$.

\medskip
\noindent\textbf{Partie I --- L'algèbre $\mathbb H$ et le
groupe $S^3$.}
\begin{enumerate}
  \item Vérifier que $\mathbb H$ est une $\R$-algèbre
        associative de \omterm{ex:b3:groups:actions}{centre} $\R$, que $\overline{pq} =
        \bar q\,\bar p$, et que $N(pq) = N(p)N(q)$
        \emph{(une voie propre~: représenter $q$ comme la
        matrice complexe $2\times2$
        $\bigl(\begin{smallmatrix}\alpha & \beta\\
        -\bar\beta & \bar\alpha\end{smallmatrix}\bigr)$, $q =
        \alpha + \beta\mathrm j$, et utiliser $\det$)}.
  \item En déduire que tout $q \neq 0$ est inversible
        ($q^{-1} = \bar q/N(q)$)~: $\mathbb H$ est un corps
        (non commutatif), et $S^3 = \{N(q) = 1\}$ est un
        groupe --- et une $3$-sous-variété compacte de
        $\R^4$ (\cref{ex:b3:submanifolds:examples}).
\end{enumerate}

\medskip
\noindent\textbf{Partie II --- Le morphisme de rotation.}
Identifier $\R^3$ aux \emph{quaternions purs} $P =
\{x\mathrm i + y\mathrm j + z\mathrm k\}$, et pour $q \in
S^3$ définir $\rho_q(v) = q\,v\,\bar q$.
\begin{enumerate}[resume]
  \item Montrer que $\rho_q$ envoie $P$ dans $P$
        \emph{(les quaternions purs sont ceux avec $\bar v =
        -v$)}, est $\R$-linéaire, préserve la norme, et que
        $\rho \colon q \mapsto \rho_q$ est un morphisme de
        groupes $S^3 \to O(3)$.
  \item Calculer le noyau~: $\rho_q = \mathrm{id}$ ssi $q$
        commute avec $\mathrm i, \mathrm j, \mathrm k$ ssi
        $q \in \R\cap S^3 = \{\pm1\}$.
  \item Écrire $q = \cos\frac\theta2 + \sin\frac\theta2\,u$
        avec $u \in P$, $N(u) = 1$ (pourquoi est-ce toujours
        possible pour $q \in S^3$~?). Montrer que $\rho_q$
        fixe $u$ et, sur le plan $u^\perp\cap P$, agit comme
        la rotation d'angle $\theta$ \emph{(calculer
        $\rho_q(w)$ pour $w \perp u$ en utilisant $uw =
        -wu$ pour des unités pures orthogonales --- prouver
        cette identité depuis la table de multiplication, ou
        depuis $uw + wu = -2\langle u, w\rangle$)}.
  \item Conclure~: $\operatorname{im}\rho \subseteq SO(3)$
        (chaque $\rho_q$ est une rotation d'axe et d'angle
        comme calculés --- déterminant $+1$ par continuité de
        $q \mapsto \det\rho_q$ sur le $S^3$ \omterm{def:b3:topology:connected}{connexe}, ou
        directement), et $\rho$ est \emph{sur} $SO(3)$~:
        toute rotation de $\R^3$ a un axe (prouver~: une
        matrice orthogonale réelle $3\times3$ de $\det = 1$
        a la valeur propre $1$ --- considérer le polynôme
        caractéristique) et est donc un certain $\rho_q$.
        Résumé~:
        \[
        SO(3) \;\cong\; S^3/\{\pm 1\} .
        \]
\end{enumerate}

\medskip
\noindent\textbf{Partie III --- $SO(3)$ comme sous-variété~;
Rodrigues.}
\begin{enumerate}[resume]
  \item Montrer que $SO(3)$ est une sous-variété compacte de
        dimension~$3$ de $M_3(\R)$ avec $T_ISO(3) =$ les
        matrices antisymétriques
        (\cref{exo:b3:submanifolds:5}~; la condition de
        déterminant sélectionne une union de composantes).
  \item Pour la matrice antisymétrique $A_u$ associée à $u
        \in \R^3$ ($A_uv = u\wedge v$, le produit vectoriel),
        prouver la \emph{formule de Rodrigues}~:
        \[
        \eu^{\theta A_u} = I + \sin\theta\,A_u + (1 -
        \cos\theta)\,A_u^2 \qquad (\norm u = 1)
        \]
        \emph{(depuis $A_u^3 = -A_u$~: scinder la série
        exponentielle le long des puissances de $A_u$)}, et
        l'identifier comme la rotation d'axe $u$ et d'angle
        $\theta$. En déduire que $\exp$ envoie les matrices
        antisymétriques \emph{sur} $SO(3)$.
  \item Relier les deux paramétrisations~: montrer que $t
        \mapsto \rho_{q(t)}$ avec $q(t) = \cos\frac t2 +
        \sin\frac t2\,u$ est un groupe à un paramètre de
        rotations dont la dérivée en $t = 0$ est $A_u$ ---
        les exponentielles quaternionique et matricielle
        racontent la même histoire à demi-vitesse et pleine
        vitesse respectivement.
\end{enumerate}

\medskip
\noindent\textbf{Partie IV --- Le double revêtement, ressenti.}
\begin{enumerate}[resume]
  \item Montrer que le chemin $q(t) = \cos\frac t2 +
        \sin\frac t2\,\mathrm k$, $t \in \intcc0{2\pi}$,
        est une boucle dans $SO(3)$ (son image $\rho_{q(t)}$
        revient à l'identité) dont le relèvement
        quaternionique n'est \emph{pas} une boucle~:
        $q(2\pi) = -q(0)$. En continuant jusqu'à $t = 4\pi$
        on ferme le relèvement. Expliquer en un court
        paragraphe ce que cela dit~: une rotation de $2\pi$
        n'est pas continûment défaisable tandis qu'une
        rotation de $4\pi$ l'est (le tour de ceinture), car
        les boucles de $SO(3)$ sont détectées dans son
        double revêtement $S^3$.
  \item En déduire aussi le dividende pratique~: composition
        des rotations = multiplication des quaternions
        ($4$ multiplications de données au lieu de $9$,
        aucune dérive hors orthogonalité) --- vérifier sur
        la composition de deux quarts de tour autour de
        $\mathrm i$ et $\mathrm j$~: calculer l'axe et
        l'angle du produit.
\end{enumerate}

\medskip
\noindent\textbf{Partie V --- La matrice explicite~:
Euler--Rodrigues.} Écrire $q = a + b\mathrm i + c\mathrm j +
d\mathrm k \in S^3$, de sorte que $a^2 + b^2 + c^2 + d^2 =
1$.
\begin{enumerate}[resume]
  \item Calculer $\rho_q(\mathrm i)$ en entier depuis la
        table de multiplication~; puis obtenir
        $\rho_q(\mathrm j)$ et $\rho_q(\mathrm k)$ par la
        substitution cyclique $\mathrm i \to \mathrm j \to
        \mathrm k \to \mathrm i$, $(b, c, d) \to (c, d, b)$
        \emph{(la justifier~: cycler $\mathrm i, \mathrm j,
        \mathrm k$ s'étend en un automorphisme de $\mathbb
        H$, car les relations définissantes sont
        cycliquement symétriques)}. Conclure que la matrice
        de $\rho_q$ dans la \omterm{def:b3:topology:basis}{base} $(\mathrm i, \mathrm j,
        \mathrm k)$ est la \emph{matrice d'Euler--Rodrigues}
        \[
        R_q = \begin{pmatrix}
        a^2 + b^2 - c^2 - d^2 & 2(bc - ad) & 2(bd + ac)\\
        2(bc + ad) & a^2 - b^2 + c^2 - d^2 & 2(cd - ab)\\
        2(bd - ac) & 2(cd + ab) & a^2 - b^2 - c^2 + d^2
        \end{pmatrix}.
        \]
  \item (Lire une rotation à rebours) Montrer que
        \[
        \operatorname{tr}R_q = 4a^2 - 1 = 1 + 2\cos\theta,
        \qquad
        \tfrac12\bigl(R_q - R_q^{\mathsf T}\bigr) =
        \sin\theta\,A_u,
        \]
        dans les notations des questions~5 et~8. En déduire
        un algorithme récupérant $\pm q$ d'une matrice de
        rotation $R$~: l'angle depuis la trace~; l'axe depuis
        la partie antisymétrique quand $0 < \theta < \pi$~;
        et, quand $\theta = \pi$, prouver et utiliser
        l'identité $R + I = 2\,uu^{\mathsf T}$.
  \item Évaluer $R_q$ pour le produit $q =
        \frac12(1 + \mathrm i + \mathrm j + \mathrm k)$ de
        la question~11~: une matrice de permutation
        apparaît. Identifier la rotation et réconcilier avec
        l'axe et l'angle trouvés à la question~11.
\end{enumerate}

\medskip
\noindent\textbf{Partie VI --- À l'intérieur de $S^3$~:
$SU(2)$, \omterm{ex:b3:groups:actions}{classes de conjugaison}, exponentielles.}
\begin{enumerate}[resume]
  \item Montrer que la \omterm{def:b3:representations:rep}{représentation} matricielle de la
        question~1 (notée $\Phi$) se restreint en un
        isomorphisme de groupes de $S^3$ sur le \emph{groupe
        unitaire spécial}
        \[
        SU(2) = \bigl\{U \in M_2(\C) : U^*U = I,\ \det U =
        1\bigr\}
        \]
        \emph{(pour la surjectivité, écrire les équations
        $U^{-1} = U^*$ et $\det U = 1$ pour une matrice
        complexe $2\times2$ générale)}.
  \item Montrer que la partie réelle est un invariant de
        conjugaison sur $S^3$ --- $\operatorname{Re}(pq\bar
        p) = \operatorname{Re}q$ pour tout $p \in S^3$ ---
        et, réciproquement, que deux quaternions unitaires de
        même partie réelle sont conjugués dans $S^3$
        \emph{(se ramener à déplacer un axe pur unitaire sur
        un autre, ce que la Partie~II fournit)}. Décrire les
        \omterm{ex:b3:groups:actions}{classes de conjugaison} de $S^3$ géométriquement~;
        traduire dans $SU(2)$ (ensembles de niveau de la
        trace)~; et projeter par $\rho$~: deux rotations
        sont conjuguées dans $SO(3)$ si et seulement si elles
        ont le même angle $\theta \in \intcc0\pi$.
  \item Définir $\exp$ sur $\mathbb H$ par la série
        exponentielle~; vérifier la convergence absolue, via
        $\abs{pq} = \abs p\,\abs q$ pour $\abs q =
        \sqrt{N(q)}$. Montrer, pour un pur unitaire $u$ et
        $\theta \in \R$,
        \[
        \exp(\theta u) = \cos\theta + \sin\theta\,u ,
        \]
        en déduire que $\exp$ envoie l'hyperplan $P$ sur
        $S^3$, et vérifier que $\rho_{\exp(su)} =
        \eu^{2sA_u}$~: le phénomène de demi-angle de la
        question~9 encore.
  \item Pour des quaternions purs $v, w$ prouver la règle de
        produit $vw = -\langle v, w\rangle + v\wedge w$,
        d'où l'identité de \omterm{def:b3:groups:derived}{commutateur} $vw - wv =
        2\,v\wedge w$~; prouver aussi $[A_v, A_w] =
        A_{v\wedge w}$ pour les matrices de la question~8.
        Conclure que la dérivée de $\rho$ en $1$ le long des
        courbes $t \mapsto \exp(tv)$ est l'isomorphisme
        linéaire $v \mapsto 2A_v$ de $P$ sur les matrices
        antisymétriques, et qu'il transporte le \omterm{def:b3:groups:derived}{commutateur}
        quaternionique vers le \omterm{def:b3:groups:derived}{commutateur} matriciel.
\end{enumerate}

\medskip
\noindent\textbf{Partie VII --- Structure globale.}
\begin{enumerate}[resume]
  \item (Pas de section \omterm{def:b3:topology:continuity}{continue}) Supposer $s \colon SO(3)
        \to S^3$ \omterm{def:b3:topology:continuity}{continue} avec $\rho \circ s =
        \operatorname{id}$. Pour la boucle $R(t) =
        \rho_{q(t)}$ de la question~10, poser
        $\varepsilon(t) = s(R(t))\,q(t)^{-1}$ pour $t \in
        \intcc0{2\pi}$. Montrer que $\varepsilon$ est
        \omterm{def:b3:topology:continuity}{continue} à valeurs dans $\{\pm1\}$, et en dériver une
        contradiction~: il n'y a pas de choix global continu
        d'un quaternion unitaire représentant chaque
        rotation.
  \item (Le modèle de la boule) Soit $\bar B \subseteq \R^3$
        la boule fermée de rayon $\pi$ et $E(v) = \eu^{A_v}$,
        avec $E(0) = I$. Montrer que $E$ envoie $\bar B$ sur
        $SO(3)$, est injective sur la boule ouverte, et sur
        la sphère frontière identifie exactement les
        antipodes~: $E(\pi u) = E(-\pi u) = 2uu^{\mathsf T}
        - I$, sans autres coïncidences. Ainsi $SO(3)$ est la
        boule avec points antipodaux du bord collés ---
        l'espace projectif $\mathbb{RP}^3$ --- et un diamètre
        devient la boucle non contractile de la question~10.
  \item Montrer que $\rho_p\rho_q\rho_p^{-1} =
        \rho_{pq\bar p}$~; que les involutions de $SO(3)$
        (les $R \neq I$ avec $R^2 = I$) sont exactement les
        demi-tours $\rho_w$ avec $w$ un quaternion pur
        unitaire~; et que le \omterm{ex:b3:groups:actions}{centre} de $SO(3)$ est trivial.
  \item Montrer que toute rotation est un produit de deux
        demi-tours~: pour $q = \cos\frac\theta2 +
        \sin\frac\theta2\,u$, choisir un pur unitaire $w
        \perp u$, vérifier que $w' = qw$ est encore un
        quaternion pur unitaire, et vérifier $\rho_q =
        \rho_{w'}\rho_w$. Où se situent les deux axes, et
        quel angle font-ils~?
  \item Conclure le résumé topologique~: $SO(3)$ est compact
        et \omterm{def:b3:topology:pathconnected}{connexe par arcs} (donner deux preuves~: image
        \omterm{def:b3:topology:continuity}{continue} de $S^3$ sous $\rho$~; image de $\exp$),
        tandis que $O(3)$ a exactement deux \omterm{def:b3:topology:components}{composantes
        connexes}, chacune homéomorphe à $SO(3)$.

  \item (Une composition, trois façons) Soient $R_1$ la
        rotation d'angle $\frac\pi2$ autour de l'axe des $z$
        et $R_2$ la rotation d'angle $\frac\pi2$ autour de
        l'axe des $x$. Calculer l'axe et l'angle de
        $R_2R_1$~: (i) en multipliant les deux matrices
        $3\times3$ et en utilisant trace/partie
        antisymétrique (Partie~V)~; (ii) en multipliant les
        quaternions unitaires correspondants $q_2q_1$.
        Vérifier que les deux réponses coïncident~: angle
        $\frac{2\pi}3$, axe $\frac1{\sqrt3}(1, -1, 1)$.
  \item (La transformée de Cayley) Pour $K$ antisymétrique,
        montrer que $I + K$ est inversible et
        \[
        C(K) = (I - K)(I + K)^{-1} \in SO(n),
        \]
        avec $-1$ jamais valeur propre de $C(K)$~; montrer
        que $K \mapsto C(K)$ est une bijection des matrices
        antisymétriques sur $\{R \in SO(n) : -1 \notin
        \operatorname{Sp}R\}$, d'inverse $R \mapsto (I -
        R)(I + R)^{-1}$. (Une carte rationnelle de $SO(n)$,
        compagne de l'$\exp$ transcendant de
        l'\cref{exo:b3:submanifolds:11}.)
\end{enumerate}
\end{problem}
